\subsection{SPACES WITH OPERATORS}

Let X be a topological space and let G be a topological group. G is said to operate continuously on X if the following conditions are satisfied:

\begin{enumerate}
\def\labelenumi{\arabic{enumi})}
\item
  X has G as a group of operators; in other words X is endowed with an external law of composition \((s, x) \to s.x\) for which G is the set of operators, and which is such that \(s.(t.x) = (st).x\) and e.x = x for all \(s, t \in G\) and all \(x \in X\) .
\item
  The mapping \((s, x) \to s.x\) of \(\mathbf{G} \times \mathbf{X}\) into \(\mathbf{X}\) is continuous.
\end{enumerate}

LEMMA 1. If a topological group G operates continuously on a topological space X, then for each \(s \in \mathbf{G}\) the mapping \(x \to s.x\) is a homeomorphism of X onto X.

For this mapping is a continuous bijection whose inverse \(x \rightarrow s^{-1}.x\) is also continuous.

We recall that for each \(x \in X\) the set G.x of transforms s.x of x by the elements s of G is called the orbit of x (with respect to the group of operators G), and that the set of all \(s \in G\) such that s.x = x is a subgroup of G called the stabilizer of x. The relation \(\mathrm{R}(x, y) : ``y \text{ belongs to the orbit of } x''\) is an equivalence relation on X, called the equivalence relation defined by G; the equivalence classes with respect to this relation are the orbits of the points of X. The topological space X/R is called the orbit space of X (with respect to G), or the quotient space of X by the group G, and is denoted by X/G; and the topology of X/G is said to be the quotient of the topology of X by G.

LEMMA 2. If a topological group G operates continuously on a topological space X, then the equivalence relation R defined by G is open.

For the saturation with respect to \(\mathbf{R}\) of an open subset \(\mathbf{U}\) of \(\mathbf{X}\) is the set \(\bigcup_{s\in G}s.\mathbf{U},\) and each \(s.\mathbf{U}\) is open by Lemma 1.

Examples. 1) Let H be a subgroup of a topological group G. H operates continuously on G by the external law \((s, x) \rightarrow sx\) . H also operates continuously on G by the external law \((x, s) \rightarrow sxs^{-1}\) .

\begin{enumerate}
\def\labelenumi{\arabic{enumi})}
\setcounter{enumi}{1}
\item
  *If K is a topological division ring (§ 6, no. 7), the multiplicative group K* operates continuously on K by the external law \((s, x) \to sx\) . *
\item
  Let G be a topological group, X a topological space. Then the mapping \((s, x) \to x\) of \(G \times X\) into X is an external law of composition on X, and G operates continuously on X with respect to this law; G is then said to operate trivially on X.
\end{enumerate}

Remark. Instead of saying that a topological group G operates continuously on a topological space X, it is often said that G operates continuously on the left on X. When the topological group \(G^{0}\) opposite to G operates continuously on X, we say that G operates continuously on the right on X. It comes to the same thing to say that X has a continuous external law of composition \((s, x) \to s.x\) with G as set of operators, such that \(s.(t.x) = (ts).x\) and e.x = x. Such a law is often written on the right: \((s, x) \to x.s\) (whence the terminology), and we have then \((x.t).s = x.(ts)\) . If G operates continuously on the right on X by the law \((s, x) \to x.s\) , then also G operates continuously on the left on X according to the external law \((s, x) \to x.s^{-1}\) , by virtue of axiom \((\mathrm{GT}_{\mathrm{II}})\) .

Let X (resp. X') be a set with a group of operators G (resp. G') and let \(f: \mathbf{G} \to \mathbf{G}'\) be a homomorphism and \(g: \mathbf{X} \to \mathbf{X}'\) a mapping. \(f\) and \(g\) are said to be compatible if \(g(s.x) = f(s).g(x)\) for all \(s \in \mathbf{G}\) and all \(x \in \mathbf{X}\) . If \(\mathbf{X}''\) is a third set with a group of operators G'', if \(f': \mathbf{G}' \to \mathbf{G}''\) is a homomorphism, \(g': \mathbf{X}' \to \mathbf{X}''\) a mapping, and if \(f'\) and \(g'\) are compatible, then \(f' \circ f\) and \(g' \circ g\) are compatible. When X, X' are topological spaces and G, G' are topological groups operating continuously on X, X' respectively, \((f, g)\) is said to be a morphism of the space with operators X into the space with operators X', provided that \(f\) and \(g\) are continuous and compatible. Passing to the quotients, \(g\) then induces a continuous mapping \(\mathbf{X}/\mathbf{G} \to \mathbf{X}'/\mathbf{G}'\) (Chapter I, § 3, no. 4, Corollary to Proposition 6).

Let G be a topological group operating continuously on a topological space X, and let \(\varphi\) be the canonical mapping of X onto the orbit space X/G. Let A be any subset of X, and let \(A'\) be the subspace of X which is the saturation of A with respect to the equivalence relation R defined by G (thus \(A'\) is the union of the orbits of points of A, and is said to be the saturation of A with respect to G). G operates continuously on \(A'\) by the restriction of \((s, x) \to s.x\) to \(G \times A'\) . Moreover, since R is open (Lemma 2) and \(A'\) is saturated, it follows from Proposition 4 of Chapter I, § 5, no. 2, and from the relation \(\varphi(A) = \varphi(A')\) , that:

PROPOSITION 10. The canonical bijection of the subspace \(\varphi(A)\) of \(X/G\) onto the orbit space \(A'/G\) is a homeomorphism.

Now let S be an equivalence relation on X such that, for each \(s \in G\) , the mapping \(x \to s.x\) is compatible with S {[}in other words, such that the relation \(x \equiv y \pmod{S}\) implies \(s.x \equiv s.y \pmod{S}\) {]}; for the sake of brevity we shall say that the relation S is compatible with the group G. If \(\psi\) is the canonical mapping of X onto X/S, and if \(s.\psi(x)\) denotes the class mod S of s.x, then X/S has G as groups of operators with respect to the external law \((s, \psi(x)) \to s.\psi(x) = \psi(s.x)\) . Moreover:

PROPOSITION 11. If the equivalence relation S on X is open and compatible with G, then G operates continuously on X/S.

Since the relation of equality on G and the relation S on X are open, it is enough to show that the mapping \((s,x)\to s.\psi(x)=\psi(s.x)\) of \(G\times X\) into X/S is continuous (Chapter I, § 5, no. 3, Corollary to Proposition 8); but this follows from the continuity of \(\psi\) and of the mapping

\[
(s, x) \rightarrow s. x.
\]

Remark. Let \(G'\) be another topological group operating continuously on X, and suppose that \(s.(s'x) = s'(s.x)\) for all \(s \in G\) , \(s' \in G'\) and \(x \in X\) . Then the equivalence relation \(S\) defined by \(G'\) is compatible with \(G\) , and since \(S\) is open (Lemma 2) it follows that \(G\) operates continuously on \(X/G'\) . Similarly, \(G'\) operates continuously on \(X/G\) . In these circumstances the two operations of the two groups \(G\) and \(G'\) on \(X\) are said to commute.

